
\section{Conclusion}
\label{sec:conclusion}
We presented \ours, a visual multi-agent self-verification framework that eliminates the need for ground-truth annotations by decomposing model responses into atomic visual claims and rewarding their consistency with the source image.
% The central contribution, Visual Claim Regulation, addresses two failure modes that undermine claim-level verification in practice---reward hacking through claim repetition and the generation of image-agnostic claims---and proves necessary for stable, non-collapsing training dynamics.
The central contribution, Visual Claim Regulation, addresses two practical failure modes in claim-level verification: reward hacking through repetitive claims and the generation of image-agnostic claims. Mitigating these issues is essential for maintaining stable and non-collapsing training dynamics.
Across 12 hallucination, VQA, and mixed benchmarks, \ours consistently outperforms competitive baselines, and scales reliably from 3B to 32B parameters.
We hope this work encourages broader adoption of structured, annotation-free self-verification as a practical path toward more grounded and trustworthy vision-language models.

% In this paper, we presented \textbf{\ours}, a multi-agent reinforced self-check 
% framework for mitigating hallucination in Vision-Language Models without relying on 
% any human annotation or ground-truth image descriptions. Inspired by the MARCH 
% paradigm for LLM factuality, we adapted the information-asymmetric Proposer-Checker 
% verification pipeline to the visual domain, where the Checker grounds each atomic 
% claim directly in the original image rather than in retrieved text documents. This 
% design breaks the self-confirmation bias inherent in single-model verification, 
% providing a more reliable and objective reward signal for GRPO training.

% To address challenges specific to the visual setting, we introduced three 
% complementary components: semantic deduplication to prevent reward inflation from 
% redundant claims, a claim quantity reward to encourage informative descriptions, and 
% an image-grounded quality penalty to suppress trivial claims answerable without 
% genuine visual attention. Together, these components stabilize training and ensure 
% that the reward signal faithfully reflects factual consistency with the image content.

% Extensive experiments on 13 benchmarks demonstrate that \ours consistently 
% outperforms both annotation-free and annotation-based baselines across diverse 
% evaluation dimensions, including hallucination detection, visual perception, and 
% image captioning. These results establish claim-level visual verification as an 
% effective and scalable paradigm for self-supervised hallucination mitigation in VLMs.

% \paragraph{Limitations and Future Work.}
% Despite its strong empirical performance, \ours has several limitations worth 
% noting. First, the Proposer-Checker pipeline introduces additional inference overhead 
% compared to single-model reward computation, as each training sample requires multiple 
% forward passes through the VLM. Reducing this cost through more efficient verification 
% strategies remains an important direction. Second, the quality of the reward signal 
% depends on the model's ability to answer visual questions accurately; for early-stage 
% training or weaker base models, the Checker may itself produce unreliable reference 
% answers, potentially introducing noise into the reward. Third, our current framework 
% focuses on image captioning tasks; extending \ours to other modalities such as 
% video understanding or interleaved image-text generation is a natural but non-trivial 
% future direction. Finally, while our quality penalty based on blurred image testing 
% effectively suppresses image-agnostic claims, more principled approaches to measuring 
% visual grounding quality---such as region-level attribution or attention-based 
% verification---may further improve the precision of the reward signal. We leave these 
% directions for future work.